\section{Riepilogo}
Nella seguente sezione viene riportato un riepilogo quantitativo dello stato dei requisiti individuati, suddivisi in base al loro grando di importanza.

\subsection{Requisiti obbligatori individuati}
\begin{table}[htbp]
	\centering
	\caption{Riepilogo requisiti obbligatori}
	\label{tab:riepilogo-obbligatori}
	\renewcommand{\arraystretch}{1.3}
	\small
	\begin{tabularx}{0.8\textwidth}{|p{4.5cm}|X|}
		\hline
		\rowcolor[gray]{0.88}
		\textbf{Tipologia Requisito} & \textbf{Numero Requisiti Obbligatori} \\
		\hline
		Funzionali & 79 \\ \hline
		Di Qualit\`a & 11 \\ \hline
		Di Vincolo & 9 \\ \hline
		Prestazionali & 3 \\ \hline
		\textbf{Totale Obbligatori} & \textbf{102} \\ \hline
	\end{tabularx}
\end{table}

\subsection{Requisiti desiderabili individuati}
\begin{table}[htbp]
	\centering
	\caption{Riepilogo requisiti desiderabili}
	\label{tab:riepilogo-desiderabili}
	\renewcommand{\arraystretch}{1.3}
	\small
	\begin{tabularx}{0.8\textwidth}{|p{4.5cm}|X|}
		\hline
		\rowcolor[gray]{0.88}
		\textbf{Tipologia Requisito} & \textbf{Numero Requisiti Obbligatori} \\
		\hline
		Funzionali & 13 \\ \hline
		Di Qualit\`a & 1 \\ \hline
		Di Vincolo & 2 \\ \hline
		Prestazionali & 0 \\ \hline
		\textbf{Totale Obbligatori} & \textbf{16} \\ \hline
	\end{tabularx}
\end{table}

\subsection{Requisiti opzionali individuati}
\begin{table}[htbp]
	\centering
	\caption{Riepilogo requisiti opzionali}
	\label{tab:riepilogo-opzionali}
	\renewcommand{\arraystretch}{1.3}
	\small
	\begin{tabularx}{0.8\textwidth}{|p{4.5cm}|X|}
		\hline
		\rowcolor[gray]{0.88}
		\textbf{Tipologia Requisito} & \textbf{Numero Requisiti Obbligatori} \\
		\hline
		Funzionali & 18 \\ \hline
		Di Qualit\`a & 0 \\ \hline
		Di Vincolo & 0 \\ \hline
		Prestazionali & 0 \\ \hline
		\textbf{Totale Obbligatori} & \textbf{18} \\ \hline
	\end{tabularx}
\end{table}

\newpage
\subsection{Riepilogo generale}
\begin{table}[htbp]
	\centering
	\caption{Riepilogo generale dei requisiti}
	\label{tab:riepilogo-generale}
	\renewcommand{\arraystretch}{1.3}
	\small
	\begin{tabularx}{0.8\textwidth}{|p{4.5cm}|X|}
		\hline
		\rowcolor[gray]{0.88}
		\textbf{Tipologia Requisito} & \textbf{Numero Requisiti Obbligatori} \\
		\hline
		Funzionali & 109 \\ \hline
		Di Qualit\`a & 12 \\ \hline
		Di Vincolo & 11 \\ \hline
		Prestazionali & 3 \\ \hline
		\textbf{Totale Obbligatori} & \textbf{135} \\ \hline
	\end{tabularx}
\end{table}